\documentclass[11pt]{article}

\usepackage[margin=1in]{geometry}
\usepackage{amsmath}
\usepackage{xcolor}
\usepackage{listings}
\usepackage{booktabs}
\usepackage{enumitem}
\usepackage[colorlinks=true,linkcolor=blue,urlcolor=blue]{hyperref}

\definecolor{codebg}{rgb}{0.96,0.96,0.96}
\definecolor{codekw}{rgb}{0.13,0.29,0.53}
\definecolor{codecm}{rgb}{0.40,0.50,0.40}
\lstset{
  basicstyle=\ttfamily\footnotesize,
  backgroundcolor=\color{codebg},
  keywordstyle=\color{codekw}\bfseries,
  commentstyle=\color{codecm}\itshape,
  breaklines=true,
  frame=single,
  framerule=0pt,
  showstringspaces=false,
  columns=fullflexible,
}

\title{\textbf{Terrain Engine} \\ \large A Real-Time OpenGL Terrain, Skybox, and Water Scene}
\author{Melvyn Qu}
\date{\today}

\begin{document}
\maketitle

\section{Overview}

Terrain Engine is a real-time 3D scene written in C++ with the OpenGL~3.3 core
profile. The viewer spawns above an island and flies freely through a world made of
three cooperating subsystems: a height-map terrain with detail multi-texturing, a
manually mapped skybox, and an animated, reflective water plane. The application is a
single flat render loop (no scene graph) so that the order and the GL state of each
pass are fully explicit.

\section{Implementation Details}

\subsection{Terrain}
The terrain mesh is generated at start-up from \texttt{heightmap.bmp}. Each grayscale
pixel $(i,j)$ becomes a vertex
\[
x = (i - \tfrac{W}{2})\,s_{xz}, \qquad
y = \tfrac{h(i,j)}{255}\,s_{y} + y_{0}, \qquad
z = (j - \tfrac{H}{2})\,s_{xz},
\]
and adjacent samples are stitched into two triangles per grid cell using an index
buffer. The color texture is sampled with the natural grid UVs, while a second
\emph{detail} texture is tiled (\texttt{uDetailTiling}) to add high-frequency surface
variation. The two textures are combined in the fragment shader with the
\emph{signed-add} rule, the modern GLSL equivalent of the fixed-function
\texttt{GL\_ADD\_SIGNED} mode recommended in the brief:
\begin{lstlisting}[language=C]
vec3 combined = clamp(baseColor + detailColor - 0.5, 0.0, 1.0);
\end{lstlisting}
Subtracting $0.5$ keeps the average brightness unchanged and avoids the colour
clamping artefacts that \texttt{GL\_ADD} or \texttt{GL\_MODULATE} would introduce.

\subsection{Skybox}
The sky is a single high-resolution \texttt{GL\_TEXTURE\_CUBE\_MAP}. It is loaded from one
horizontal-cross PNG ($4\times3$ cells); the six faces are sliced straight out of the
decoded image using \texttt{GL\_UNPACK\_ROW\_LENGTH}, so no external image tooling is
needed. The cube is rendered by sampling the cubemap with the per-vertex direction
$\texttt{mat3(model)}\cdot\texttt{aPos}$. A pleasant consequence is that the water
\emph{reflection} comes for free: the reflection pass already scales the cube by
$(1,-1,1)$, which flips the sampling direction in $Y$ and produces a mirrored sky with no
extra code. Wrapping is \texttt{GL\_CLAMP\_TO\_EDGE} on all three axes to avoid edge
seams. (An earlier version mapped five separate 2D textures onto the cube faces with
manual per-face UV rotations; the cubemap replaced that, both sharpening the sky and
removing the orientation bookkeeping.)

\subsection{Water}
The bottom plane samples a 128$\times$128 wave texture with \texttt{GL\_REPEAT} and a
tiling factor well above $1.0$ so the low-resolution image is not stretched. Animation
is achieved by scrolling the UVs every frame by an accumulated \texttt{uWaveShift}
(advanced by $\text{speed}\times\Delta t$ so the motion is frame-rate independent):
\begin{lstlisting}[language=C]
vec2 uv = vec2(TexCoords.x * uWaterTiling + uWaveShift,
               TexCoords.y * uWaterTiling + 0.5 * uWaveShift);
\end{lstlisting}
The surface is alpha-blended over the scene and tinted slightly toward a sea-blue. A
distance-based \texttt{smoothstep} fade makes the water more transparent near the
camera (revealing the reflection) and opaque toward the horizon, which hides the far
edge of the plane.

\subsection{Reflections}
Two reflections are rendered. The \emph{sky} reflection re-draws the skybox cube scaled
by $(1,-1,1)$ below the water line with depth writes disabled. The \emph{terrain}
reflection multiplies the terrain model matrix by a mirror transform about the water
plane, $T(0,\,2y_{w},\,0)\cdot S(1,-1,1)$. Because the negative scale reverses triangle
winding, the front-face convention is temporarily switched to \texttt{GL\_CW} for that
draw and restored afterwards. Each terrain pass is bounded by a user clip plane fed
through \texttt{gl\_ClipDistance[0]}: the reflection keeps only geometry below the water
surface, and the real pass keeps only geometry above it.

\subsection{Camera and asset paths}
Movement uses a standard Euler fly-camera (\texttt{camera.h}) driven by WASD plus
Space/Shift, mouse-look, and scroll-to-zoom. As required by the brief, motion is
\emph{smooth}: instead of teleporting the position each frame, the camera keeps a
velocity vector that exponentially approaches the target speed while a key is held
(gradual acceleration to a cruising speed) and decays back to zero when released
(gradual deceleration). The exponential approach $v \mathrel{+}= (v_{\text{target}}-v)
\,(1-e^{-r\,\Delta t})$ makes the feel frame-rate independent.
All asset and shader paths are resolved from a compile-time \texttt{PROJECT\_ROOT}
string injected by CMake, so the executable loads its resources correctly regardless of
the current working directory.

\section{Challenges \& Solutions}

\begin{description}[leftmargin=0pt,style=nextline]
\item[Machine-specific hard-coded paths]
The project originally embedded machine-specific absolute paths (rooted under a
\texttt{computer\_graphics/} home directory) for every shader and texture, and one of
them pointed at a directory that no longer existed, so asset loading failed at runtime. All paths were rewritten to be relative to the
CMake-provided \texttt{PROJECT\_ROOT}, and the build configuration was repointed at the
in-repository \texttt{external/} dependencies.

\item[Broken / missing build dependencies]
The build referenced a GLFW \texttt{.dylib} that was no longer present (only the static
\texttt{libglfw3.a} remained) and a FreeType library that was not installed. The fix was
to link the static GLFW archive together with the required macOS frameworks
(\texttt{Cocoa}, \texttt{IOKit}, \texttt{CoreVideo}, \texttt{OpenGL}), and to install the
remaining toolchain (CMake, and FreeType while the text overlay still existed).

\item[Skybox seams and orientation]
With the default clamp mode the cube edges showed visible seams, and several face images
were rotated relative to the cube. Switching to \texttt{GL\_CLAMP\_TO\_EDGE} removed the
seams, and per-face UV-rotation uniforms aligned every image.

\item[Reflection winding]
Mirroring the terrain with a negative scale inverted the triangle winding, which caused
the reflected terrain to be culled. Temporarily setting \texttt{glFrontFace(GL\_CW)}
around the mirrored draw resolved it.

\item[HUD text orientation (documented, then removed)]
An on-screen FreeType HUD displayed the camera position, but the glyphs rendered flipped
because FreeType stores bitmaps top-row-first while the texture sampling assumed the
opposite. After several iterations on the vertical flip, the HUD was intentionally
removed from the final build to keep the scene clean and to drop the FreeType
dependency entirely. This is a deliberate scope decision, not an unresolved bug.
\end{description}

\paragraph{Note on orientation.} Because cubemap faces follow a fixed convention,
swapping in a new sky is now a matter of replacing one cross PNG; if a particular set is
authored with a different face order, only the cross-cell mapping in
\texttt{loadCubemapCross} needs adjusting.

\section{Environment Setup}

\begin{tabular}{ll}
\toprule
Operating system & macOS (Darwin 25.4, Apple Silicon) \\
Compiler / standard & Apple Clang, C++14 \\
Build system & CMake ($\geq 3.5$) \\
Windowing / input & GLFW 3.4 (static \texttt{libglfw3.a}) \\
GL loader & GLAD (OpenGL 3.3 core) \\
Math library & GLM \\
Image loading & stb\_image \\
\bottomrule
\end{tabular}

\medskip
\noindent GLAD, GLFW, GLM, and stb\_image are vendored under \texttt{external/}; only
CMake must be installed to build. The application is built with
\texttt{cmake -B build \&\& cmake --build build} and runs as \texttt{./terrainEngine}.

\section{Bonus \& References}

\subsection*{Features beyond the basic requirements}
\begin{itemize}[leftmargin=1.5em]
  \item \textbf{Dual reflection} of both the sky and the terrain in the water, each
        clipped to the correct half-space with \texttt{gl\_ClipDistance}.
  \item \textbf{Distance-based water transparency} using a \texttt{smoothstep} fade, so
        the water reveals the reflection up close and hides the plane edge at range.
  \item \textbf{Frame-rate-independent} wave animation via $\Delta t$ accumulation.
  \item \textbf{Portable asset resolution} through a compile-time \texttt{PROJECT\_ROOT}
        define, allowing the binary to run from any directory.
  \item \textbf{Self-contained build}: all third-party dependencies vendored in-repo.
  \item \textbf{4$\times$ MSAA} for anti-aliased geometry edges.
  \item \textbf{High-resolution cubemap sky} sliced in-code from a single cross image.
\end{itemize}

\subsection*{Possible future improvements}
\begin{itemize}[leftmargin=1.5em]
  \item Per-vertex or normal-mapped \textbf{terrain lighting} (currently unlit albedo).
  \item \textbf{GPU-displaced} water with real normals for specular sun glint, instead of
        a scrolling texture.
  \item \textbf{Fresnel-weighted} reflection blending for a more physically convincing
        water surface.
  \item \textbf{Frustum culling / LOD} on the terrain grid for larger maps.
  \item Higher-resolution or procedurally generated sky.
\end{itemize}

\subsection*{References}
\begin{enumerate}[leftmargin=1.5em]
  \item Joey de Vries, \emph{LearnOpenGL} --- \url{https://learnopengl.com/}
        (camera, shaders, text-rendering, skybox tutorials).
  \item GLFW 3.4 documentation --- \url{https://www.glfw.org/docs/latest/}.
  \item OpenGL Mathematics (GLM) --- \url{https://github.com/g-truc/glm}.
  \item Sean Barrett, \texttt{stb\_image} --- \url{https://github.com/nothings/stb}.
  \item Course design brief, \texttt{doc/TerrainDoc.md} (skybox / water / terrain spec).
  \item ``Cloudy Skyboxes'' by Screaming Brain Studios, CC0 / public domain ---
        \url{https://opengameart.org/content/cloudy-skyboxes-0} (skybox cubemap art).
\end{enumerate}

\end{document}
